\documentclass[11pt]{article}
\usepackage[margin=1in]{geometry}
\usepackage{fontspec}
\setmainfont{Arial}
\usepackage{graphicx}
\usepackage{fvextra}
\usepackage{enumitem}
\setlist{nosep}
\setlength{\parindent}{0pt}
\setlength{\parskip}{0.55em}
\begin{document}

\begin{center}\LARGE\textbf{Agreeableness evidence specimen (QY2 skill-forward validation)}\end{center}
\section*{1 · QA inputs}
Input bindings: two answered Probes copied locally into this View's own

input/QA-probes/ folder.

\begin{Verbatim}[fontsize=\small]
input/QA-probes/1-canonical-definition.md → canonical trait meaning
input/QA-probes/2-observable-signal.md    → observable signal + bound source count
\end{Verbatim}
This View imports two answered QA Probes, both copied (not moved, not edited)

from the QBt1 specimen's own input folder for this skill-validation exercise.

\begin{quote}\footnotesize\textbf{Evidence card.} Card two answered QA Probes: QI1 · QA input · Bindings: views/QY2-for-view/input/QA-probes/1-canonical-definition.md, views/QY2-for-view/input/QA-probes/2-observable-signal.md. Status: answered. Role: supply the material organized by this View.\end{quote}
\section*{2 · View body}
Readable body: the topic-specific material a person should understand before

opening any Card.

\subsection*{2.1 · Trait meaning}
Agreeableness is an interpersonal trait expressed through warmth and cooperation.

\begin{quote}\footnotesize\textbf{Evidence card.} Card interpersonal trait: EC1 · Literature · Finding: agreeableness is treated as an interpersonal construct. Bearing: supports. Boundary: trait description only. Binding: views/QY2-for-view/input/QA-probes/1-canonical-definition.md. Freshness: current. Used by: QY2-Display1.\end{quote}
The literature-side construct frame resolves through the real bibliography (John et al., 1999).

\subsection*{2.2 · Observable project signal}
Review text records patient descriptions of patient-facing behavior, and the

current View binds 3 source Probe records into one evidence judgment, so this

View treats review-derived scores as patient-perceived signals rather than an

error-free latent trait measure.

\begin{quote}\footnotesize\textbf{Evidence card.} Card 3 source Probe records: EC2 · Value · Finding: the project observes descriptions and evaluations of behavior across 3 bound source records. Bearing: qualifies. Boundary: observed signal, not direct latent trait. Binding: views/QY2-for-view/input/QA-probes/2-observable-signal.md. Freshness: current. Used by: QY2-Display1.\end{quote}
The body is not a hidden evidence ledger. It is the human-readable synthesis

produced from the two inputs above; the Citation and Value Cards let a reader

inspect why each exact phrase is present.

\section*{3 · Displays}
Display outputs: one rendered table selects and compares both View-body

findings.

QY2-Display1 is the agreeableness evidence table. Its current raster is embedded

below; opening QY2-Display1 shows the same preview first, before its View-body

bindings, files, and independent acceptance state.

\begin{center}\includegraphics[width=0.94\linewidth]{assets/display-1.png}\end{center}
\begin{quote}\footnotesize\textbf{Evidence card.} Card QY2-Display1: D1 · Display · kind: table · reader job: compare the canonical trait definition and the observable project signal. Bindings: EC1, EC2. Binding: views/QY2-for-view/output/QY2-Display1-agreeableness-evidence-table/output.md. Status: rendered. Acceptance: waiting.\end{quote}
preview.pdf is the printable inspection surface for the same table; it sits

alongside preview.png in the same Display folder.

\section*{4 · Consumers}
Consumer bindings: one local Results-section stand-in inside the QY2 sandbox

plans to use QY2-Display1.

The Results section plans to consume QY2-Display1 for its construct-interpretation passage.

\begin{quote}\footnotesize\textbf{Evidence card.} Card Results section: C1 · Consumer · Target Page: S-Results · Binding: consumers/S-Results.md. Uses: QY2-Display1. Placement: Results / construct interpretation. Status: planned, blocked on View and Display acceptance.\end{quote}
\section{References}\begin{itemize}
\item John, Oliver P.; Srivastava, Sanjay (1999). The Big-Five Trait Taxonomy: History, Measurement, and Theoretical Perspectives. Handbook of Personality: Theory and Research
\end{itemize}
\end{document}
